% Generated by scripts/assemble.py; corrections belong in source/corrections.json.
\examhead{60}{天津大学 2026 年数学分析试卷}{59}
% p060-b19; source PDF page 60
\blockanchor{p060-b19}
\question{1}(10 分) 设 $f$ 是 $\mathbb{R}^2$ 上的连续函数, 且 $f$ 在原点沿任意方向的方向导数存在. 问: $f$ 在原点是否一定可微? 若可微, 给出证明, 若不可微, 给出反例.\par
% p060-b20; source PDF page 60
\blockanchor{p060-b20}
\question{2}(15 分) 证明: $f(x) = \sin \left( \frac{1}{x^2} \right)$ 在 $(0,1]$ 上非一致连续, 在 $[1, +\infty)$ 上一致连续.\par
% p060-b21; source PDF page 60
\par\Needspace{6\baselineskip}
\blockanchor{p060-b21}
\question{3}(15 分) 已知函数\par
\nopagebreak[4]
% p060-b22; source PDF page 60
\blockanchor{p060-b22}
\displayformula{f (x) = \left\{ \begin{array}{l l} A, & 0 \leq x \leq \pi ; \\ - A, & - \pi \leq x <   0. \end{array} \right.}
% p060-b23; source PDF page 60
\blockanchor{p060-b23}
其中 $A \neq 0$ 为常数, 求该函数在区间 $[- \pi, \pi]$ 的傅里叶级数, 并证明 $\sum_{n=1}^{\infty} \frac{1}{(2n-1)^2} = \frac{\pi^2}{8}$.\par
% p060-b24; source PDF page 60
\blockanchor{p060-b24}
\question{4}(10 分) 设函数 $f(x)$ 在 $[0, +\infty)$ 的任何有限闭区间上可积, 令 $L(s) = \int_{0}^{+\infty} e^{-sx} f(x) \, \mathrm{d}x$, 已知 $L(s_0)$ 收敛, 证明: $L(s)$ 在 $s \geq s_0 > 0$ 时收敛.\par
% p060-b25; source PDF page 60
\blockanchor{p060-b25}
\question{5}(15 分) 设  $f(x) \in C[0,1]$ ，且  $\int_{0}^{1} f(x)e^{nx} dx = 0$  对所有  $n = 0, 1, 2, \cdots$  成立，证明： $f(x) \equiv 0$ .\par
% p061-b00; source PDF page 61
\blockanchor{p061-b00}
\question{6}(15 分) 设  $\boldsymbol{F} = (F^{1}, F^{2}, \cdots, F^{n})$ :  $R^{n} \to R^{n}$  在原点可微, 且  $\boldsymbol{F}(\boldsymbol{0}) = \boldsymbol{0}$ . 已知  $\sum_{i,j=1}^{n} \left( \frac{\partial F^{i}(\boldsymbol{0})}{\partial x_{j}} \right)^{2} < 1$ , 证明: 存在  $\delta > 0$ , 当  $|x| \leq \delta$  时, 有  $|\boldsymbol{F}(\boldsymbol{x})| \leq \delta$ .\par
% p061-b01; source PDF page 61
\blockanchor{p061-b01}
\question{7}(15 分) 设函数 $f, g: \mathbb{R} \to \mathbb{R}$ 连续可微, 且 $f(0) = 0$, $f'(0) \neq 0$, 考虑方程 $f(x) = tg(x)$\par
% p061-b02; source PDF page 61
\blockanchor{p061-b02}
(1) 证明: 存在 $\delta > 0$, 使得 $x(t)$ 是方程在 $t \in (-\delta, \delta)$ 上满足 $x(0) = 0$ 的唯一解.\par
% p061-b03; source PDF page 61
\blockanchor{p061-b03}
(2) 求 $x(t)$ 在 $t = 0$ 处的一阶泰勒展开式.\par
% p061-b04; source PDF page 61
\blockanchor{p061-b04}
\question{8}(20 分) 设 $B$ 是 $\mathbb{R}^3$ 中的单位开球, 函数 $u \in C^2(B) \cap C(\overline{B})$ 满足 (i) $u$ 是 $B$ 上的非零函数, (ii) $u|_{\partial B} = 0$, (iii) 存在 $\lambda \in \mathbb{R}$, 使得在 $B$ 上有 $\Delta u = \lambda u$.\par
% p061-b05; source PDF page 61
\blockanchor{p061-b05}
(1) 证明: $\int_{B} |\nabla u|^{2} \, \mathrm{d}x \, \mathrm{d}y \, \mathrm{d}z + \lambda \int_{B} u^{2} \, \mathrm{d}x \, \mathrm{d}y \, \mathrm{d}z = 0$.\par
% p061-b06; source PDF page 61
\blockanchor{p061-b06}
(2) 证明: $\lambda < 0$.\par
% p061-b07; source PDF page 61
\blockanchor{p061-b07}
\question{9}(20 分) 设 $\{b_n\}$ 是严格单调递增趋于 $+\infty$ 的数列, 且 $\lim_{n \to \infty} \frac{a_n - a_{n-1}}{b_n - b_{n-1}} = A (A \in \mathbb{R})$.\par
% p061-b08; source PDF page 61
\blockanchor{p061-b08}
(1) 证明: $\lim_{n\to \infty}\frac{a_n}{b_n} = A.$\par
% p061-b09; source PDF page 61
\blockanchor{p061-b09}
(2) 计算  $\lim_{n\to\infty}\frac{1^{k}+2^{k}+\cdots+n^{k}}{n^{k+1}}$ ，其中 k 为正整数.\par
% p061-b10; source PDF page 61
\blockanchor{p061-b10}
\question{10}(15 分) 证明: 自然数 $e$ 是无理数.\par
